\section{Projective reduction to an orbit--slope Bessel theorem}
\label{sec:tpc42-projective-orbit}

The residue obstruction does not make the low projective complexity of
the physical masks useless.  It tells us exactly what that complexity can
and cannot reduce.  We now pass from the full Hilbert coefficient to one
declared arithmetic form on each TPC--36 layer.

\subsection{Exact one-layer norm identity}

Fix one of the fixed number of orbit residue cells produced in TPC--36 and
replace \(\cJ\) by its intersection with that cell; we keep the notation
\(\cJ\).  On this fixed cell TPC--36 gives the exact representation
\begin{equation}
 \mathfrak A_{\alpha,\gamma}^{\rm act}(j)
 =\int_\Omega
 p_\omega(\alpha)q_\omega(\gamma)
 \eta_\omega(j/J)\,d\nu(\omega),
 \label{eq:tpc42-projective-representation}
\end{equation}
and, for every fixed smoothness order needed below,
\begin{equation}
 \int_\Omega
 \|p_\omega\|_\infty
 \|q_\omega\|_\infty
 \|\eta_\omega\|_{C^B}\,d|\nu|(\omega)
 \ll_B X^{o(1)}.
 \label{eq:tpc42-projective-mass}
\end{equation}
The factors may be normalized to have sup norm at most one, with their
sizes absorbed into \(d|\nu|\).  Put
\begin{equation}
 \Gamma_{\omega,\alpha}
 :=\gamma_\alpha^{(1)}p_\omega(\alpha).
 \label{eq:tpc42-layer-source}
\end{equation}
The left output of one layer at alias \(k\) is
\begin{equation}
 [\mathscr T_{\omega,k}^L]_{\gamma,j}
 =q_\omega(\gamma)\eta_\omega(j/J)
  \sum_\alpha\Gamma_{\omega,\alpha}b_{\alpha,j,k}.
 \label{eq:tpc42-one-layer-output}
\end{equation}

\begin{proposition}[Exact projective layer reduction]
\label{prop:tpc42-projective-layer-reduction}
For every projective layer,
\begin{align}
 \|\mathscr T_{\omega,k}^L\|_2^2
 &=\|q_\omega\|_2^2
   \sum_{j\in\cJ}
   |\eta_\omega(j/J)|^2
   \left|
    \sum_\alpha\Gamma_{\omega,\alpha}b_{\alpha,j,k}
   \right|^2,
 \label{eq:tpc42-layer-fixed-k-identity}\\
 \sum_{k\in\cI_B}\|\mathscr T_{\omega,k}^L\|_2^2
 &=\|q_\omega\|_2^2
   \sum_{j\in\cJ}|\eta_\omega(j/J)|^2
   \sum_{k\in\cI_B}
   \left|
    \sum_\alpha\Gamma_{\omega,\alpha}b_{\alpha,j,k}
   \right|^2.
 \label{eq:tpc42-layer-trace-identity}
\end{align}
The corresponding layer diagonals are
\begin{align}
 \cD_{\omega,k}^L
 &=\|q_\omega\|_2^2
   \sum_j|\eta_\omega(j/J)|^2
   \sum_\alpha
   |\Gamma_{\omega,\alpha}|^2|b_{\alpha,j,k}|^2,
 \label{eq:tpc42-layer-fixed-diagonal}\\
 \cD_{\omega,\rm tr}^L
 &=\sum_{k\in\cI_B}\cD_{\omega,k}^L.
 \label{eq:tpc42-layer-trace-diagonal}
\end{align}
\end{proposition}

\begin{proof}
In \eqref{eq:tpc42-one-layer-output}, the dependence on \(\gamma\) is the
single factor \(q_\omega(\gamma)\).  Square, sum first over \(\gamma\),
and then over \(j\).  This proves
\eqref{eq:tpc42-layer-fixed-k-identity}; summing over \(k\) proves the
trace identity.  Deleting the unequal-row products inside the same
squares gives \eqref{eq:tpc42-layer-fixed-diagonal} and
\eqref{eq:tpc42-layer-trace-diagonal}.
\end{proof}

The proposition shows that Hilbert dualization introduces no mysterious
new coordinate after projective separation.  It leaves the explicit
orbit--slope form
\begin{equation}
 \sum_{j\in\cJ}w_\omega(j)
 \left|
  \sum_\alpha\Gamma_{\omega,\alpha}
  \mu(m_\alpha j+h+kM)
  A_k(m_\alpha j)
 \right|^2,
 \qquad
 w_\omega(j)=|\eta_\omega(j/J)|^2.
 \label{eq:tpc42-orbit-slope-form}
\end{equation}
For fixed \(j\), all rows in this sum occupy distinct target residues
modulo \(M\).  It is therefore a transversal row estimate, not a variance
inside one arithmetic progression.

The same conclusion follows directly by duality.  If
\(z=(z_{\gamma,j})\) is a unit vector, then a layer contributes
\begin{equation}
 \sum_j\eta_\omega(j/J)
 \left(\sum_\gamma q_\omega(\gamma)
       \overline{z_{\gamma,j}}\right)
 \sum_\alpha\Gamma_{\omega,\alpha}b_{\alpha,j,k},
 \label{eq:tpc42-layer-duality}
\end{equation}
and
\begin{equation}
 \sum_j\left|
  \sum_\gamma q_\omega(\gamma)
  \overline{z_{\gamma,j}}
 \right|^2
 \le\|q_\omega\|_2^2\|z\|_2^2.
 \label{eq:tpc42-dual-orbit-bound}
\end{equation}
Conversely, every orbit vector in the direction of \(q_\omega\) is
attained.  Thus \eqref{eq:tpc42-layer-fixed-k-identity} is the exact
duality norm, not a Cauchy majorant.

\subsection{The coefficient-specific continuation}

The direct equality route would follow from the following layer-stable
estimate:
\begin{align}
 &\sum_j|\eta_\omega(j/J)|^2
 \left|
  \sum_\alpha\Gamma_{\omega,\alpha}b_{\alpha,j,0}
 \right|^2
 \notag\\
 &\qquad\ll X^{o(1)}K_B
 \sum_j|\eta_\omega(j/J)|^2
 \sum_\alpha
 |\Gamma_{\omega,\alpha}|^2|b_{\alpha,j,0}|^2.
 \label{eq:tpc42-orbit-slope-equality-target}
\end{align}
The cleaner full-trace route asks for
\begin{align}
 &\sum_{j,k}|\eta_\omega(j/J)|^2
 \left|
  \sum_\alpha\Gamma_{\omega,\alpha}b_{\alpha,j,k}
 \right|^2
 \notag\\
 &\qquad\ll X^{o(1)}
 \sum_{j,k}|\eta_\omega(j/J)|^2
 \sum_\alpha
 |\Gamma_{\omega,\alpha}|^2|b_{\alpha,j,k}|^2.
 \label{eq:tpc42-orbit-slope-trace-target}
\end{align}
Here and below the quantifier ranges only over the projective layers
actually generated by TPC--36, their fixed residue cells, and the two
physical row systems.  It does not range over arbitrary bounded source
coefficients.

\begin{proposition}[Uniform layer theorem implies the endpoint bounds]
\label{prop:tpc42-layer-to-physical}
If \eqref{eq:tpc42-orbit-slope-equality-target} holds uniformly on all
actual normalized layers and residue cells, then
\begin{equation}
 \|\mathscr T_0^L\|^2+\|\mathscr T_0^R\|^2
 \ll X^{o(1)}K_BQ^2J.
 \label{eq:tpc42-layer-to-equality-endpoint}
\end{equation}
If \eqref{eq:tpc42-orbit-slope-trace-target} holds uniformly instead,
then
\begin{equation}
 \sum_{k\in\cI_B}
 \left(\|\mathscr T_k^L\|^2+\|\mathscr T_k^R\|^2\right)
 \ll X^{o(1)}Q^2JK_B.
 \label{eq:tpc42-layer-to-trace-endpoint}
\end{equation}
These are absolute endpoint bounds; no comparison with the atomic
diagonal of the signed reassembled kernel is asserted.
\end{proposition}

\begin{proof}
Normalize \(p_\omega,q_\omega,\eta_\omega\) to have sup norm at most one
and absorb their former sizes into the signed layer measure.  Then its
total variation is \(P=X^{o(1)}\) by
\eqref{eq:tpc42-projective-mass}.  Moreover,
\begin{equation}
 \|q_\omega\|_2^2\ll X^{o(1)}Q,
 \qquad
 \sum_\alpha|\Gamma_{\omega,\alpha}|^2
 \ll X^{o(1)}Q,
 \qquad
 |\cJ|\ll J,
 \label{eq:tpc42-layer-endpoint-ledgers}
\end{equation}
by \eqref{eq:tpc42-source-l2}, \eqref{eq:tpc42-row-counts}, and the
normalization.  The terminal logarithms cost only \(X^{o(1)}\).

The equality-layer hypothesis and
\cref{prop:tpc42-projective-layer-reduction} therefore give
\begin{equation}
 \|\mathscr T_{\omega,0}^L\|^2
 \ll X^{o(1)}K_BQ^2J
 \label{eq:tpc42-one-layer-equality-absolute}
\end{equation}
uniformly.  Minkowski's integral inequality costs \(P^2=X^{o(1)}\)
after squaring and proves the left part of
\eqref{eq:tpc42-layer-to-equality-endpoint}.  Interchanging the row
systems gives the right part.

For the trace hypothesis, work in
\(\ell^2(\cI_B)\otimes\Hh\).  The same three ledgers, now with
\(|\cI_B|=K_B\), give
\begin{equation}
 \sum_k\|\mathscr T_{\omega,k}^L\|^2
 \ll X^{o(1)}Q^2JK_B.
 \label{eq:tpc42-one-layer-trace-absolute}
\end{equation}
Minkowski in this product Hilbert space again costs only \(P^2\), with no
additional alias factor.  The fixed residue-cell reassembly is absorbed
in \(X^{o(1)}\), completing the proof.
\end{proof}

Normalization before applying the uniform layer theorem is necessary.  A
bound proved separately after renormalizing each formal layer, followed
by an absolute count of the
\(O(D_0)\) Ferrers layers, would reintroduce a polynomial loss.  TPC--36
supplies small signed projective mass, not a small raw layer count.

\begin{remark}[What projective reduction does not prove]
\label{rem:tpc42-projective-boundary}
The constant kernel used in
\cref{thm:tpc42-scalar-blind-countermodel} is itself a projective layer of
norm one.  Hence \eqref{eq:tpc42-projective-mass} alone implies neither
\eqref{eq:tpc42-orbit-slope-equality-target} nor
\eqref{eq:tpc42-orbit-slope-trace-target}.  The new input must exploit the
actual arithmetic direction \(\Gamma_{\omega,\alpha}\), including its
source M\"obius sign, jointly with the target M\"obius values along
\(m_\alpha j+h+kM\).  Establishing that restricted Bessel theorem would
be genuine progress beyond scalar progression variance; its statement
does not claim a parity breach or a prime-pair asymptotic.
\end{remark}
